\documentclass[12pt,a4paper]{article}
\usepackage[utf8]{inputenc}
\usepackage[T5]{fontenc}
\usepackage[vietnamese]{babel}
\usepackage[margin=1.5cm,top=1.6cm,bottom=1.5cm]{geometry}
\usepackage{amsmath,amssymb,array,tabularx,enumitem,fancyhdr,tikz,xcolor}
\usetikzlibrary{calc,angles,quotes,arrows.meta,babel}
\definecolor{ink}{RGB}{30,55,70}
\newcommand{\headerQuote}{%
  \ifcase\value{page}\relax
    \or Math gives us hope that every problem has a solution% Trang 1
    \or Math + You = Champion% Trang 2
    \or Where Frustration Becomes Inspiration% Trang 3
    \or Turning Confusion into Clarity% Trang 4
    \or Don't be an X, be a Y of mathematics% Trang 5
    \or Don't get mad. Learn to add% Trang 6
    \or Math gives us hope that every problem has a solution% Trang 7
    \or Math + You = Champion% Trang 8
    \or Where Frustration Becomes Inspiration% Trang 9
    \or Turning Confusion into Clarity% Trang 10
    \or Don't be an X, be a Y of mathematics% Trang 11
    \or Don't get mad. Learn to add% Trang 12
    \else Math gives us hope that every problem has a solution%
  \fi
}
\pagestyle{fancy}\fancyhf{}
\lhead{\small\textit{\headerQuote}}
\rhead{\small\textit{Người biên soạn: Dương Gia Bảo}}
\rfoot{\small Trang \thepage}
\renewcommand{\headrulewidth}{0.4pt}
\setlength{\headheight}{18pt}
\setlength{\parindent}{0pt}
\setlength{\parskip}{2pt}
\setlength{\tabcolsep}{5pt}
\renewcommand{\arraystretch}{1.16}
\newcommand{\parttitle}[1]{\par\vspace{3pt}{\large\bfseries\color{ink}#1}\par\vspace{2pt}}
\newcommand{\four}[4]{\par\noindent\begin{tabularx}{\linewidth}{@{}*{4}{>{\raggedright\arraybackslash}X}@{}}\textbf{A.} #1&\textbf{B.} #2&\textbf{C.} #3&\textbf{D.} #4\end{tabularx}}
\newcommand{\two}[4]{\par\noindent\begin{tabularx}{\linewidth}{@{}*{2}{>{\raggedright\arraybackslash}X}@{}}\textbf{A.} #1&\textbf{B.} #2\\[3pt]\textbf{C.} #3&\textbf{D.} #4\end{tabularx}}
\newcommand{\answerbox}{\hfill\textit{Đáp số:}\ \fbox{\rule{0pt}{13pt}\hspace{2.2cm}}}
\setlist[enumerate,1]{label=\textbf{Câu \arabic*.},leftmargin=1.5cm,labelsep=.18cm,itemsep=3.5pt,topsep=2pt,parsep=0pt}
\setlist[enumerate,2]{label=\alph*),leftmargin=.65cm,itemsep=2.5pt,topsep=2pt,parsep=0pt}

\begin{document}
\noindent
\begin{minipage}[c]{0.38\textwidth}
    \centering
    \textbf{LÀM TOÁN ĐỂ HẠNH PHÚC}\\[-2pt]
    \rule{3.2cm}{0.6pt}\\[4pt]
    \textbf{ĐỀ CHÍNH THỨC}
\end{minipage}\hfill
\begin{minipage}[c]{0.35\textwidth}
    \centering
    \textbf{ĐỀ ÔN THI GK1}\\
    \textbf{NĂM HỌC 2026 -- 2027}\\[2pt]
    \textbf{ĐỀ SỐ 02}\\
    \textbf{Môn: TOÁN -- Khối lớp 10}\\
    \textit{\small Thời gian làm bài: 90 phút\\(không kể thời gian phát đề)}
\end{minipage}

\vspace{0.2cm}
\noindent
\begin{minipage}[b]{0.78\textwidth}
    \textbf{Họ và tên:}\dotfill\textbf{SBD:}\makebox[3.5cm]{\dotfill}
\end{minipage}\hfill
\begin{minipage}[b]{0.20\textwidth}
    \hfill
    \setlength{\fboxrule}{0.8pt}%
    \fbox{\begin{minipage}{2.3cm}
        \centering
        \textbf{\small Mã đề thi}\\
        \textbf{0002}
    \end{minipage}}
\end{minipage}
\vspace{0.2cm}

\parttitle{PHẦN I. Trắc nghiệm nhiều lựa chọn}
\textit{Trả lời câu 1 đến câu 12. Mỗi câu chỉ chọn một phương án.}
\begin{enumerate}
\item Trong các câu sau, câu nào là một mệnh đề toán học?
\four{Số $13$ là một số nguyên tố.}{Hôm nay trời đẹp quá!}{Bạn có thích học môn Toán không?}{Hãy giải phương trình này ngay!}
\item Cho mệnh đề chứa biến $P(x)\colon\text{ ``}x^2 - 4 = 0\text{''}$ với $x\in\mathbb{R}$. Khẳng định nào sau đây là một mệnh đề đúng?
\four{$P(1)$.}{$P(2)$.}{$P(0)$.}{$P(3)$.}
\item Cho tập hợp $X = \{a; b; c\}$. Khẳng định nào sau đây đúng?
\four{$a \subset X$.}{$\{b\} \in X$.}{$\{a; b\} \subset X$.}{$d \in X$.}
\item Số tập con của tập hợp $S = \{1; 2; 3; 4\}$ là
\four{$4$.}{$8$.}{$12$.}{$16$.}
\item Bất phương trình nào sau đây là bất phương trình bậc nhất hai ẩn $x, y$?
\four{$3x - 5y > 7$.}{$2x^2 + y \leq 3$.}{$xy - 2 \geq 0$.}{$x - \dfrac{2}{y} < 1$.}
\item Cặp số $(x; y)$ nào sau đây là một nghiệm của bất phương trình $2x - y + 1 > 0$?
\four{$(0; 2)$.}{$(1; 1)$.}{$(-1; 0)$.}{$(1; 4)$.}
\item Hệ bất phương trình nào sau đây là hệ bất phương trình bậc nhất hai ẩn $x, y$?
\two{$\begin{cases} x - 2y \leq 1 \\ x^2 + y > 0 \end{cases}$}{$\begin{cases} 2x + y < 5 \\ x - y^2 \geq 2 \end{cases}$}{$\begin{cases} 2x - y \leq 3 \\ x + 3y > 4 \end{cases}$}{$\begin{cases} xy \leq 2 \\ x + y > 1 \end{cases}$}
\item Điểm nào sau đây thuộc miền nghiệm của hệ bất phương trình $\begin{cases} x \geq 0 \\ y \geq 0 \\ x + y \leq 5 \\ 2x + y \geq 4 \end{cases}$?
\four{$(0; 1)$.}{$(3; 3)$.}{$(-1; 5)$.}{$(2; 2)$.}
\item Cho góc $\alpha$ thỏa mãn $0^\circ < \alpha < 90^\circ$. Khẳng định nào sau đây đúng?
\four{$\sin\alpha > 0$.}{$\cos\alpha < 0$.}{$\tan\alpha < 0$.}{$\cot\alpha < 0$.}
\item Với mọi góc $\alpha$ thỏa mãn $0^\circ \leq \alpha \leq 180^\circ$, khẳng định nào sau đây đúng?
\two{$\sin(180^\circ - \alpha) = -\sin\alpha$.}{$\cos(180^\circ - \alpha) = -\cos\alpha$.}{$\tan(180^\circ - \alpha) = \tan\alpha$ ($\alpha \neq 90^\circ$).}{$\cot(180^\circ - \alpha) = \cot\alpha$ ($0^\circ < \alpha < 180^\circ$).}
\item Cho tam giác $ABC$ có các cạnh $BC = a, CA = b, AB = c$ và bán kính đường tròn ngoại tiếp là $R$. Khẳng định nào sau đây đúng?
\four{$\dfrac{a}{\sin A} = R$.}{$\dfrac{a}{\cos A} = 2R$.}{$\dfrac{a}{\sin A} = 2R$.}{$\dfrac{b}{\sin B} = \dfrac{R}{2}$.}
\item Cho tam giác $ABC$ có nửa chu vi $p = \dfrac{a+b+c}{2}$ và bán kính đường tròn nội tiếp $r$. Diện tích $S$ của tam giác $ABC$ được tính theo công thức nào sau đây?
\four{$S = \dfrac{p}{r}$.}{$S = 2pr$.}{$S = \dfrac{r}{p}$.}{$S = pr$.}
\end{enumerate}

\par\vspace{8pt}
\parttitle{PHẦN II. Trắc nghiệm đúng/sai}
\textit{Với mỗi ý a), b), c), d) trong từng câu, hãy xác định đúng hoặc sai.}
\begin{enumerate}
\item Cho hai tập hợp $A = [-3; 5)$ và $B = (0; 7]$.
\begin{enumerate}
\item $-3 \in A$.
\item $5 \in A$.
\item $A \cap B = (0; 5)$.
\item $A \setminus B = [-3; 0]$.
\end{enumerate}
\item Một xưởng may gia công cần may hai loại áo: áo sơ mi (gọi số lượng là $x$) và áo khoác (gọi số lượng là $y$), với $x, y$ là các số nguyên không âm. Mỗi chiếc áo sơ mi cần $1{,}5$ m vải và $2$ giờ công may. Mỗi chiếc áo khoác cần $2{,}5$ m vải và $3$ giờ công may. Xưởng có tối đa $120$ m vải và $130$ giờ công may.
\begin{enumerate}
\item Tổng lượng vải cần dùng để may hai loại áo trên là $1{,}5x + 2{,}5y$ (m).
\item Kế hoạch may $30$ chiếc áo sơ mi và $30$ chiếc áo khoác vượt quá tổng lượng vải hiện có của xưởng.
\item Các phương án may khả thi thỏa mãn hệ bất phương trình
$\begin{cases} 3x + 5y \leq 240, \\ 2x + 3y \leq 130, \\ x, y \in \mathbb{Z},\ x \geq 0,\ y \geq 0. \end{cases}$
\item Nếu xưởng quyết định may đúng $35$ chiếc áo sơ mi thì có thể may tối đa $21$ chiếc áo khoác.
\end{enumerate}
\end{enumerate}

\vspace{7pt}
\parttitle{PHẦN III. Trả lời ngắn}
\textit{Ghi kết quả vào ô đáp số. Chỉ làm tròn khi câu hỏi yêu cầu.}
\begin{enumerate}
\item Cho hai tập hợp $A = \{x \in \mathbb{Z} \mid -5 < 2x - 1 \leq 9\}$ và $B = \{x \in \mathbb{Z} \mid |x| \leq 3\}$. Tập hợp $A \cap B$ có bao nhiêu phần tử?
\answerbox
\item Cho góc $\alpha$ thỏa mãn $90^\circ < \alpha < 180^\circ$ và $\sin\alpha = \dfrac{3}{5}$. Tính giá trị của biểu thức
\[
M = \frac{5\cos\alpha + 2\tan\alpha}{1 + \cos\alpha}.
\]
\answerbox
\item \begin{minipage}[t]{0.48\linewidth}
Để đo khoảng cách giữa hai vị trí $A$ và $B$ bị ngăn cách bởi một hồ nước, người ta chọn một điểm mốc $C$ trên bờ sao cho từ $C$ có thể quan sát trực tiếp cả $A$ và $B$. Người đo xác định được khoảng cách $CA = 60\text{ m}$, $CB = 80\text{ m}$ và góc ngắm $\widehat{ACB} = 60^\circ$. Khoảng cách giữa hai điểm $A$ và $B$ bằng bao nhiêu mét? \textit{(Làm tròn kết quả đến hàng phần mười)}.
\par\vspace{4pt}\answerbox
\end{minipage}\hfill%
\begin{minipage}[t]{0.50\linewidth}
\vspace{-0.4\baselineskip}
\centering
\begin{tikzpicture}[scale=1.1, font=\normalsize, >=Stealth, line join=round, line cap=round]
    % Mặt hồ minh hoạ
    \fill[blue!8, rounded corners=15pt] (0.8,0.3) -- (3.8,0.5) -- (4.2,2.8) -- (1.5,2.6) -- cycle;
    \node[blue!70!black, font=\small\bfseries] at (2.5,1.5) {Hồ nước};
    
    % Các điểm
    \coordinate (C) at (0.2,0.4);
    \coordinate (A) at (1.5,3.2);
    \coordinate (B) at (4.5,0.4);
    
    % Nối các cạnh
    \draw[line width=1.4pt, ink] (C) -- node[above left=2pt, font=\normalsize\bfseries, black] {$60\text{ m}$} (A);
    \draw[line width=1.4pt, ink] (C) -- node[below=3pt, font=\normalsize\bfseries, black] {$80\text{ m}$} (B);
    \draw[dashed, line width=1.5pt, red!80!black] (A) -- node[above right=1pt, font=\normalsize\bfseries, red!80!black] {$?$} (B);
    
    % Góc C
    \pic[draw, line width=1.1pt, angle radius=8mm, angle eccentricity=1.4, "$60^\circ$", font=\normalsize\bfseries] {angle=B--C--A};
    
    % Điểm đỉnh
    \fill (C) circle (3pt) node[below left=2pt, font=\Large\bfseries] {$C$};
    \fill (A) circle (3pt) node[above=2pt, font=\Large\bfseries] {$A$};
    \fill (B) circle (3pt) node[below right=2pt, font=\Large\bfseries] {$B$};
\end{tikzpicture}
\end{minipage}
\item Gọi $D$ là tập hợp tất cả các điểm $(x; y)$ thỏa mãn hệ bất phương trình
\[
x \geq 0,\quad y \geq 0,\quad x + y \leq 5,\quad x + 2y \leq 8.
\]
Tìm giá trị lớn nhất của tham số thực $m$ để bất phương trình $x + (m - 2)y \leq 10$ nghiệm đúng với \textbf{mọi} điểm $(x; y) \in D$.
\par\answerbox
\end{enumerate}

\par\vspace{8pt}
\parttitle{PHẦN IV. Tự luận}
\textit{Trình bày lập luận và các bước tính trên giấy làm bài.}
\begin{enumerate}
\item Cho hai tập hợp $A = [-4; 2)$ và $B = [-1; 6]$.
Xác định các tập hợp $A \cap B$, $A \cup B$, $A \setminus B$ và biểu diễn tập hợp $A \setminus B$ trên trục số.
\item Cho góc $\alpha$ thỏa mãn $0^\circ < \alpha < 90^\circ$ và biểu thức
\[
Q = \frac{\sin^2\alpha - \cos^2\alpha}{\sin\alpha\cos\alpha + \cos^2\alpha}.
\]
\begin{enumerate}
\item Chứng minh rằng $Q = \tan\alpha - 1$.
\item Tính giá trị của biểu thức $Q$ khi $\cos\alpha = \dfrac{5}{13}$.
\end{enumerate}
\item Từ hai trạm quan sát $A$ và $B$ trên một bờ biển thẳng cách nhau $500\text{ m}$, các kiểm soát viên theo dõi một con tàu du lịch $C$ đang neo đậu ngoài biển. Người ta đo được các góc $\widehat{CAB} = 75^\circ$ và $\widehat{CBA} = 60^\circ$ như hình vẽ bên dưới.
\begin{enumerate}
\item Tính khoảng cách từ trạm $A$ đến con tàu $C$ (đoạn $AC$), \textit{làm tròn đến hàng đơn vị theo mét}.
\item Tính khoảng cách ngắn nhất từ con tàu $C$ đến bờ biển nối $A$ và $B$ (chiều cao $CH$ hạ từ $C$ xuống đường thẳng $AB$), \textit{làm tròn đến hàng đơn vị theo mét}.
\end{enumerate}
\begin{center}
\begin{tikzpicture}[scale=1.2, line join=round, line cap=round, font=\normalsize]
    % Đường bờ biển
    \draw[line width=1.4pt, teal!80!black] (-0.8,0) -- (6.0,0) node[right=2pt, font=\normalsize\bfseries, black] {Bờ biển};
    
    % Tọa độ các điểm
    \coordinate (A) at (0,0);
    \coordinate (B) at (4.5,0);
    % Góc A = 75, B = 60 => C = 45. AC = 4.5 * sin(60)/sin(45) ~ 5.51. H nằm trên AB
    % Tọa độ C: x = AC * cos(75) = 5.51 * 0.2588 = 1.426; y = AC * sin(75) = 5.51 * 0.9659 = 5.32
    \coordinate (C) at (1.426, 3.8);
    \coordinate (H) at (1.426, 0);
    
    % Vùng biển
    \fill[cyan!8] (-0.8,0) -- (6.0,0) -- (6.0,4.2) -- (-0.8,4.2) -- cycle;
    \node[blue!70!black, font=\normalsize\bfseries] at (4.2,2.5) {Biển};
    
    % Các đường ngắm
    \draw[line width=1.5pt, ink] (A) -- node[above left=1pt, font=\normalsize\bfseries, black] {$AC$} (C);
    \draw[line width=1.5pt, ink] (B) -- (C);
    \draw[dashed, line width=1.2pt, red!80!black] (C) -- node[right=2pt, font=\normalsize\bfseries, red!80!black] {$h$} (H);
    \draw[line width=1.5pt, ink] (A) -- node[below=3pt, font=\normalsize\bfseries, black] {$500\text{ m}$} (B);
    
    % Ký hiệu góc vuông và góc ngắm
    \pic[draw, line width=1.0pt, angle radius=4mm] {right angle=C--H--B};
    \pic[draw, line width=1.1pt, angle radius=7mm, angle eccentricity=1.5, "$75^\circ$", font=\normalsize\bfseries] {angle=B--A--C};
    \pic[draw, line width=1.1pt, angle radius=8mm, angle eccentricity=1.45, "$60^\circ$", font=\normalsize\bfseries] {angle=C--B--A};
    
    % Các đỉnh
    \fill (A) circle (3pt) node[below left=2pt, font=\Large\bfseries] {$A$};
    \fill (B) circle (3pt) node[below right=2pt, font=\Large\bfseries] {$B$};
    \fill (C) circle (3pt) node[above=2pt, font=\Large\bfseries] {$C$};
    \fill (H) circle (2.5pt) node[below=2pt, font=\normalsize\bfseries] {$H$};
\end{tikzpicture}
\end{center}
\item Một cơ sở thủ công mỹ nghệ sản xuất hai loại giỏ mây: loại I và loại II theo từng kiện sản phẩm.
Để sản xuất một kiện loại I cần $3\text{ kg}$ mây và $2$ giờ công lao động, đem lại tiền lãi $400\,000$ đồng.
Để sản xuất một kiện loại II cần $2\text{ kg}$ mây và $4$ giờ công lao động, đem lại tiền lãi $500\,000$ đồng.
Mỗi ngày, cơ sở có tối đa $36\text{ kg}$ mây và $48$ giờ công lao động. Giả sử sản phẩm làm ra đều được tiêu thụ hết.
\begin{enumerate}
\item Gọi $x, y$ lần lượt là số kiện giỏ mây loại I và loại II mà cơ sở sản xuất trong một ngày ($x, y \in \mathbb{N}$). Lập hệ bất phương trình biểu thị các điều kiện của bài toán và viết biểu thức tiền lãi $T(x; y)$ thu được mỗi ngày.
\item Cơ sở nên sản xuất bao nhiêu kiện mỗi loại để tiền lãi thu được mỗi ngày là lớn nhất? Tính số tiền lãi lớn nhất đó.
\end{enumerate}
\end{enumerate}
\vfill
\begin{center}\textbf{--- HẾT ---}\end{center}

\newpage
\begin{center}
{\Large\bfseries BẢNG ĐÁP ÁN TỔNG HỢP}\\[5pt]
{\bfseries ÔN TẬP GIỮA KỲ I -- ĐỀ SỐ 02}
\end{center}
\parttitle{PHẦN I. Trắc nghiệm nhiều lựa chọn}
\begin{center}
\begin{tabularx}{\linewidth}{|>{\bfseries}l|*{12}{>{\centering\arraybackslash}X|}}\hline
Câu&1&2&3&4&5&6&7&8&9&10&11&12\\\hline
Đáp án&A&B&C&D&A&B&C&D&A&B&C&D\\\hline
\end{tabularx}
\end{center}

\parttitle{PHẦN II. Trắc nghiệm đúng/sai}
\begin{center}
\begin{tabular}{|c|c|c|c|c|}\hline
\textbf{Câu}&\textbf{a)}&\textbf{b)}&\textbf{c)}&\textbf{d)}\\\hline
1&Đúng&Sai&Đúng&Đúng\\\hline
2&Đúng&Sai&Đúng&Sai\\\hline
\end{tabular}
\end{center}

\parttitle{PHẦN III. Trả lời ngắn}
\begin{center}
\begin{tabularx}{\linewidth}{|>{\bfseries}l|*{4}{>{\centering\arraybackslash}X|}}\hline
Câu&1&2&3&4\\\hline
Đáp số&$5$&$-\dfrac{55}{2}$&$72{,}1\text{ m}$&$4{,}5$\\\hline
\end{tabularx}
\end{center}

\parttitle{PHẦN IV. Tự luận -- kết quả cần đạt}
\begin{tabularx}{\linewidth}{|c|>{\raggedright\arraybackslash}X|}\hline
\textbf{Câu}&\textbf{Kết quả tóm tắt}\\\hline
1&$A\cap B = [-1; 2)$; $A\cup B = [-4; 6]$; $A\setminus B = [-4; -1)$.\\
&Trục số của $A\setminus B$: lấy nửa khoảng $[-4; -1)$, gạch bỏ phần còn lại bên ngoài.\\[4pt]\hline
2&a) Rút gọn: $Q = \tan\alpha - 1$.\\
&b) Khi $\cos\alpha = \dfrac{5}{13}$ thì $\sin\alpha = \dfrac{12}{13} \implies \tan\alpha = \dfrac{12}{5} \implies Q = \dfrac{7}{5}$.\\[4pt]\hline
3&a) $AC = \dfrac{500\sin 60^\circ}{\sin 45^\circ} = 250\sqrt{6} \approx \boxed{612\text{ m}}$.\\
&b) Chiều cao $CH = AC\sin 75^\circ = 125(3+\sqrt{3}) \approx \boxed{592\text{ m}}$.\\[4pt]\hline
4&a) Hệ BPT: $\begin{cases} 3x + 2y \leq 36 \\ x + 2y \leq 24 \\ x, y \in \mathbb{Z},\ x \geq 0,\ y \geq 0 \end{cases}$; Tiền lãi: $T(x; y) = 400\,000x + 500\,000y$ (đồng).\\
&b) Cơ sở nên sản xuất \textbf{6 kiện loại I} và \textbf{9 kiện loại II}; tiền lãi lớn nhất đạt $\boxed{6\,900\,000\text{ đồng}}$.\\\hline
\end{tabularx}

\newpage
\begin{center}
{\Large\bfseries HÌNH VẼ MIỀN NGHIỆM HỆ BẤT PHƯƠNG TRÌNH}\\[3pt]
{\normalsize\bfseries (PHẦN III -- CÂU 4 \& PHẦN IV -- CÂU 4)}
\end{center}

\vspace{6pt}
\begin{center}
\textbf{\large 1. Phần III -- Câu 4: Miền đa giác $D$ của hệ bất phương trình}\\[6pt]
\begin{tikzpicture}[scale=1.2, font=\normalsize, >=Stealth, line join=round, line cap=round]
    % Trục toạ độ
    \draw[->, line width=1.2pt] (-1.0,0) -- (6.5,0) node[below=2pt, font=\normalsize\bfseries] {$x$};
    \draw[->, line width=1.2pt] (0,-0.8) -- (0,5.8) node[left=2pt, font=\normalsize\bfseries] {$y$};
    \node[below left=2pt, font=\normalsize\bfseries] at (0,0) {$O$};
    
    % Miền D tô màu
    \fill[blue!15] (0,0) -- (5,0) -- (2,3) -- (0,4) -- cycle;
    
    % Đường thẳng d1: x + y = 5 (kéo từ góc trên trái xuống x=5)
    \draw[line width=1.3pt, blue!80!black] (-0.5,5.5) -- (5.5,-0.5);
    \node[above right, font=\normalsize\bfseries, blue!80!black] at (3.8,1.6) {$d_1\colon x+y=5$};
    
    % Đường thẳng d2: x + 2y = 8 (kéo từ Oy qua B(2;3))
    \draw[line width=1.3pt, teal!80!black] (-0.8,4.4) -- (6.5,0.75);
    \node[above left, font=\normalsize\bfseries, teal!80!black] at (0.2,4.6) {$d_2\colon x+2y=8$};
    
    % Cạnh miền đa giác D (nét rất đậm)
    \draw[line width=2.0pt, ink] (0,0) -- (5,0) -- (2,3) -- (0,4) -- cycle;
    
    % Các đỉnh của miền D
    \fill (0,0) circle (3.5pt);
    \fill (5,0) circle (3.8pt);
    \node[below=4pt, font=\normalsize\bfseries, black] at (5,0) {$A(5;0)$};
    \fill (2,3) circle (4.0pt) node[above right=3pt, font=\normalsize\bfseries, black] {$B(2;3)$};
    \fill (0,4) circle (3.8pt) node[left=4pt, font=\normalsize\bfseries, black] {$C(0;4)$};
    
    % Đường dóng đỉnh B(2;3)
    \draw[dashed, line width=1.0pt, gray!80!black] (2,0) -- (2,3) -- (0,3);
    \node[below left=1pt, font=\normalsize\bfseries, black] at (2,0) {$2$};
    \node[left=4pt, font=\normalsize\bfseries, black] at (0,3) {$3$};
    
    % Nhãn miền D
    \node[ink, font=\Large\bfseries] at (1.5, 1.4) {Miền $D$};
\end{tikzpicture}
\end{center}

\vspace{6pt}
\hrule
\vspace{6pt}

\begin{center}
\textbf{\large 2. Phần IV -- Câu 4: Miền nghiệm tứ giác $OABC$ của bài toán tối ưu}\\[6pt]
\begin{tikzpicture}[scale=0.75, font=\normalsize, >=Stealth, line join=round, line cap=round]
    % Trục toạ độ
    \draw[->, line width=1.2pt] (-1.5,0) -- (15.0,0) node[below=2pt, font=\normalsize\bfseries] {$x$ (kiện loại I)};
    \draw[->, line width=1.2pt] (0,-1.0) -- (0,14.5) node[left=2pt, font=\normalsize\bfseries] {$y$ (kiện loại II)};
    \node[below left=2pt, font=\normalsize\bfseries] at (0,0) {$O$};
    
    % Miền nghiệm tô màu
    \fill[ink!15] (0,0) -- (12,0) -- (6,9) -- (0,12) -- cycle;
    
    % Đường 3x + 2y = 36 (đi qua A(12;0), B(6;9), (0;18))
    \draw[line width=1.3pt, blue!80!black] (3.0,13.5) -- (13.0,-1.5);
    \node[above right, font=\normalsize\bfseries, blue!80!black] at (3.5, 13.0) {$3x+2y=36$};
    
    % Đường x + 2y = 24 (đi qua C(0;12), B(6;9))
    \draw[line width=1.3pt, teal!80!black] (-1.0,12.5) -- (14.5,4.75);
    \node[above right, font=\normalsize\bfseries, teal!80!black] at (12.5, 5.8) {$x+2y=24$};
    
    % Cạnh miền nghiệm tứ giác
    \draw[line width=2.0pt, ink] (0,0) -- (12,0) -- (6,9) -- (0,12) -- cycle;
    
    % Các đỉnh
    \fill (0,0) circle (3.8pt);
    \fill (12,0) circle (4.0pt);
    \node[below right=1pt, font=\normalsize\bfseries, black] at (12,0) {$A(12;0)$};
    \fill (6,9) circle (4.2pt) node[above right=3pt, font=\normalsize\bfseries, black] {$B(6;9)$};
    \fill (0,12) circle (4.0pt) node[left=4pt, font=\normalsize\bfseries, black] {$C(0;12)$};
    
    % Đường dóng đỉnh B(6;9)
    \draw[dashed, line width=1.0pt, gray!80!black] (6,0) -- (6,9) -- (0,9);
    \node[below left=1pt, font=\normalsize\bfseries, black] at (6,0) {$6$};
    \node[left=4pt, font=\normalsize\bfseries, black] at (0,9) {$9$};
    
    % Nhãn miền nghiệm
    \node[ink, font=\Large\bfseries] at (4.0, 3.8) {Miền tứ giác $OABC$};
\end{tikzpicture}
\end{center}
\end{document}
